Decentralized exchanges (DEXs) let anyone observe a pending trade before it executes. Block producers and third-party searchers exploit this to reorder, insert or censor transactions for profit, a practice known as Maximal Extractable Value (MEV). The sandwich attack, in which an attacker trades immediately before and after a user's swap, is the most direct form of this harm: the attacker's profit is taken from the user.

Existing defences are partial. Private transaction channels hide trades from outside bots but not from the relay or block builder that carries them. RediSwap goes further by auctioning MEV and refunding the proceeds to users and liquidity providers, but it depends on a private sequencer. A compromised sequencer can read transactions and extract MEV before the auction takes place, so it is a single point of failure.

This dissertation presents \textbf{MEVRES}, a mechanism that keeps RediSwap's redistribution but removes the trusted sequencer. A user commits to a trade with a hash and encrypts it under a threshold public key whose decryption key is shared among $n$ nodes, any $t+1$ of which must cooperate to decrypt. Arbitrageurs then bid in a \textbf{second-price auction} without seeing the trade. Only after the bids are fixed is the trade decrypted, checked against its commitment, and executed inside an atomic sandwich bundle that returns the pool to its initial state. The winning arbitrageur's payment is refunded to the user and to liquidity providers.

We prove that truthful bidding is a dominant strategy for arbitrageurs, that participation is individually rational for arbitrageurs, users and liquidity providers, that an arbitrageur cannot gain by bidding under several identities, and that bundles for different transactions are independent of one another.

We implemented a prototype in which the hashing, encryption and secret sharing are executed for real, and simulated the mechanism over 1{,}000 independent runs. Under a worst-case attacker model, MEVRES reduced the average cost of trading for users from 164 to 36 basis points; when only 5\% of profitable trades are attacked, the reduction is still 52\%. An ablation shows that encryption without redistribution reaches 70 basis points, so the auction halves the cost again. Liquidity provider returns turned from a loss into a profit, and about 78\% of the auctioned MEV was returned to users and 13\% to liquidity providers. A further experiment shows how the benefit of a RediSwap-style design decays as its sequencer becomes less trustworthy.

Finally, we implemented the on-chain part of MEVRES as a Solidity contract and benchmarked it with Hyperledger Caliper. All 11{,}700 benchmark transactions succeeded. With six bidders and two trades per epoch a MEVRES trade costs about 10.7 times the gas of a plain swap, falling to about 4 times when 50 trades share an epoch.

\vspace{1em}
\noindent\textbf{Keywords:} Maximal Extractable Value, Decentralized Finance, Sandwich Attack, Threshold Encryption, Commitment Scheme, Second-Price Auction, Constant Function Market Maker, MEV Redistribution, Blockchain, Hyperledger Caliper.
